\documentclass[11pt]{article}
\usepackage[margin=1in]{geometry}\usepackage{amsmath,amssymb}\usepackage{booktabs}\usepackage{hyperref}\usepackage{enumitem}
\title{Virtualizing Decision and Context Spaces for Decision-Native Agents}\author{Anonymous}\date{}
\begin{document}\maketitle
\begin{abstract}We propose a runtime for decision-native agents that virtualizes both an unbounded decision space and an unbounded context space. A replaceable Decision Model maps state and resident options to a distribution over typed decisions, $D(s,O)\to P(O)$. Virtual Option Space uses paging and decision-preserving refinement to expose a bounded resident option set; Virtual Context Space maintains pinned, working, and cold context blocks and recovers from ContextFaults through staged paging. The design composes coverage, granularity, semantic freshness, and consistency mechanisms while remaining backend agnostic. Jev is the current prototype backend, not the definition of the abstraction. This draft reports a mechanism baseline and small Jev replay only; evaluation results are placeholders.\end{abstract}
\section{Introduction}Generative agents typically treat actions and context as prompt contents. Decision-native agents instead receive structured state and emit typed decisions. Their bottleneck is therefore runtime management of what can be considered and what is resident. We study Open World execution through dual virtualization: logical options become resident options, and logical context becomes a working set. Our claims concern the composition and interfaces, not independent novelty of each policy.
\section{Motivation}A finite decision interface cannot enumerate an unbounded option set. Likewise, retaining all historical state can dilute decision utility. Conventional RAG retrieves information for answering; here memory/RAG is a context residency policy. The runtime does not rely on cross-request prefix/KV-cache reuse, allowing aggressive context mutation; this makes no claim that Jev has no internal KV cache. We distinguish semantic freshness from cache locality.
\section{Decision Model abstraction}The backend interface is $D(s,O,C)\to P(O)$, returning probabilities and typed metadata for resident options $O$ under resident context $C$. Implementations may be Jev, a mock, an oracle, or another typed decision service. A state transition commits only validated decisions and increments revisions.
\section{Dual Virtualization}The logical option and context spaces are unbounded. Managers materialize bounded resident options and blocks, then invoke the same Decision Model. Option metadata includes identity, parent, granularity, dependencies, revision, and expansion status. Context blocks include block\_id, summary, raw\_ref, revision, dependencies, last\_access, access\_count, utility, type, size, and pinned.
\section{Virtual Option Space}PAGE/EXPAND broadens coverage by loading more same-granularity candidates. We define progressive refinement as $\operatorname{Refine}(o,s)\to\mathcal{O}'$: given a coarse option $o$ and state $s$, the runtime constructs a finer-grained option space $\mathcal{O}'$ with explicit parent identity, revision, dependencies, and budgets. Thus PAGE/EXPAND is horizontal coverage expansion, while REFINE is vertical resolution expansion. REFINE replaces a coarse option with compatible finer candidates while preserving provenance and consistency; it does not transfer final authority from the Decision Model to a proposer. Resident options use hysteresis and minimum residency to prevent churn. REVISION and INVALIDATE remove stale choices before commit.
\section{Virtual Context Space}Pinned context contains goals, hard constraints, and irreversible decisions and does not age. Working blocks are the current resident set; Cold blocks remain addressable in compressed or indexed storage. Type-specific aging protects task state, while observations and transient retrieval age faster. Utility aging combines decay with observed decision usefulness. Phase-aware replacement and working-set history balance semantic freshness against locality.
\section{PAGE, REFINE, and ContextFault}The fast path retrieves a working set and invokes $D$. Low coverage, ambiguous scores, high decision entropy, multi-hop classification, or an explicit insufficient-context decision raises ContextFault. Recovery pages summaries or clusters first, then materializes concrete blocks and retries the same logical decision. Option faults similarly trigger PAGE/EXPAND or REFINE. A refined child carries its parent ID, revision, dependency signature, and budget accounting; stale parents or exhausted depth, resident, byte, call, or risk budgets produce a RefineFault and bounded recovery. Recovery is bounded and recorded for evaluation.
\section{State consistency and recovery}Every option and block carries a revision and dependency set. Before commit, the runtime checks revisions, invalidates conflicting residents, and replays from the last consistent state when needed. Fault recovery must be idempotent; duplicate retries cannot commit two transitions.
\section{Implementation mapping}The current prototype exposes a replaceable \texttt{DecisionModel} boundary with Jev/Choice, replay, and oracle adapters. It implements mechanism-level option paging through \texttt{VirtualOptionManager}, plus deterministic context residency through \texttt{ContextResidencyManager}. These managers are currently separate replayable slices rather than a complete production scheduler. Jev is replaceable. No claim is made that the prototype is a complete production system.
\section{Evaluation plan}Results are TODO. We will first close the real Jev loop, then evaluate three converged lines: (A) decision-space virtualization, comparing fixed resident, retrieval top-$k$, hierarchical candidates, and virtual paging; (B) progressive refinement and recovery, with a controlled \texttt{COMMIT/PAGE/REFINE/FALLBACK/CLARIFY/STOP} confusion matrix; and (C) dynamic context residency, comparing static, append-only, retrieval replacement, aging, and anti-thrashing policies. A 20--100 step decision-dense workload will connect the three lines. We will vary logical options (10, 100, 1K, 10K, 100K), resident $K\in\{8,16,32\}$, and backend capability (60,70,80,90,oracle). Missing detection, page localization, recovery success, and end-to-end success are separate metrics; an oracle-known target page is only a mechanism upper bound. Metrics are RecoveryRate, coverage, decision cost, latency, thrashing, and semantic freshness. The current evidence is a mechanism baseline and small Jev replay, not an end-to-end quality result. A direct masked-diffusion proposal failure will be retained only as an appendix negative result.
\section{Related Work}MemGPT/Letta virtualize LLM memory; CMV, Pichay, RAPTOR, GraphRAG, and llm-mmu provide adjacent context hierarchy or paging ideas. ToolChain* and hierarchical tool selectors motivate structured action selection. AnyJev and Mini-Jev expose Jev-style typed decisions, while Jev-Mem, jev-memory, JevHarness, REFLEX, and Hermes Jev Skills demonstrate Jev memory, selective control, routing, or agent tooling. FUDGE, GeDi, constrained decoding, and speculative decoding motivate helper-generated candidate tokens. Our bounded claim is narrower: in the public literature and implementations searched through 2026-09-24, we did not find this exact composition of a replaceable DecisionModel runtime that virtualizes logical options as a demand-paged resident set, defines PAGE/REFINE/INVALIDATE semantics, and co-designs cache-agnostic context residency. This is a scoped statement about the searched record, not an absolute claim of global priority.
\section{Limitations and Threats}The prototype has no broad benchmark, learned utility model, or production Jev backend study. Oracle comparisons may overestimate achievable quality. Semantic freshness is task dependent; paging and aggressive mutation can thrash despite safeguards. Backend internals and cache behavior are not assumed or measured.
\section{Conclusion}Decision-native agents benefit from a runtime that virtualizes both what they may choose and what they may see. A replaceable DecisionModel, resident options, resident context, refinement, ContextFault recovery, and revision checks provide a testable mechanism baseline. The planned evaluation will determine whether the composition improves coverage and cost.
\bibliographystyle{plain}\bibliography{references}
\end{document}
